\chapter{Conclusion}
\label{cha:conclusion}

\section{Summary of Work}
\label{sec:summary}

This thesis set out to design, implement, and evaluate an automated pre-processing
pipeline for industrial pyrometer temperature data from metal forming and
heat-treatment processes. The work was motivated by a practical need: raw
pyrometer signals from press-hardening lines cannot be used directly for process
monitoring or control because they contain measurement noise, impulse spikes, and
generate data volumes that are difficult to store and transmit efficiently.

Two pipeline stages were addressed in this individual contribution  signal
denoising~(ATP-1) and data compression~(ATP-3)  while the sensor calibration
stage~(ATP-2) is addressed in a parallel thesis by Avula Ajay Kumar. Ten
denoising methods and six compression methods were implemented in Python and
evaluated on the open-source NIST additive manufacturing dataset across all ten
available layers. The five classical denoising and compression methods were
additionally validated on ten real two-colour pyrometer scans from the PODFAM
dataset, confirming that the best classical methods generalise well to real
industrial pyrometer data.

The complete pipeline is delivered as a set of fifteen modular, well-commented
Python scripts available at
\url{https://github.com/Sravya0304/pyrometer-preprocessing}.

\section{Key Findings}
\label{sec:findings}

The most important finding from the denoising evaluation is that the Kalman filter
clearly outperforms all other methods on both datasets. It achieves 89.5\,\%
noise reduction on NIST and 88.4\,\% on PODFAM, far ahead of any ML method.
This result is explained by the nature of the signal: the NIST camera signal is
dominated by Gaussian electronic noise on a slowly varying temperature trend,
which is exactly the noise type the Kalman filter was designed to handle.

The ML denoisers~(CNN, LSTM, Autoencoder, BiLSTM) achieve reductions of only
27--32\,\% on NIST. The main limiting factors are the quality of the training
target~(median-filtered signal rather than a physical thermocouple measurement)
and the small training dataset of approximately 1,652 frames from a single NIST
layer.

For compression, wavelet compression with 10\,\% Haar coefficient retention
provides the best balance of ratio and accuracy: $5.0\times$ on NIST with RMSE
36.55\,\textdegree C, satisfying the ATP-3 target. This is the only method that
meets both the ratio target and maintains a reconstruction error small enough for
process monitoring applications. PCA achieves the highest ratio~($10.7\times$)
but at greater reconstruction error.

The cross-dataset consistency of the classical methods is also noteworthy. On
PODFAM the Kalman filter achieves 88.4\,\% noise reduction in every one of the
ten files, and wavelet compression achieves approximately~$3.1\times$ ratio with
RMSE below~2\,\textdegree C in every file, all without any retuning.

\section{Answers to Research Questions}
\label{sec:rq-answers}

\textbf{Research Question~1:} \textit{How effectively can classical filtering
methods and machine learning-based denoisers reduce the noise in industrial
pyrometer signals, and which method provides the best noise reduction while
preserving genuine temperature peaks?}

The Kalman filter provides the best noise reduction of all ten methods:
89.5\,\% on NIST and 88.4\,\% on PODFAM. The LSTM achieves the best ML result
on NIST~(31.8\,\%). The median filter is recommended for industrial deployment
because it combines reliable spike suppression with peak preservation on both
datasets without requiring training data. ATP-1 is satisfied --- all methods
reduce noise relative to the raw baseline, and the best classical method achieves
a reduction above 89\,\%.

\textbf{Research Question~2:} \textit{What compression ratio is achievable for
calibrated pyrometer time-series data, and what is the trade-off between
compression ratio and reconstruction accuracy?}

Wavelet compression with 10\,\% Haar coefficient retention achieves the best
classical trade-off: $5.0\times$ ratio with RMSE~36.55\,\textdegree C on NIST,
satisfying the ATP-3 target of~$CR > 4\times$. PCA achieves the highest
ratio~($10.7\times$) at greater reconstruction error~(58.1\,\textdegree C).
Downsampling and SVD do not satisfy the ratio target. ATP-3 is satisfied by
the wavelet method on NIST.

\section{Limitations}
\label{sec:limitations-ch6}

\textbf{Simulated thermocouple reference.} The NIST dataset contains no physical
thermocouple. The simulated reference~($\alpha = 0.08$ low-pass filter) is an
approximation, not a real measurement. Reported denoising RMSE values are
therefore relative comparisons between methods rather than absolute accuracy
estimates.

\textbf{Limited ML training data.} Approximately~1,652 frames from NIST Layer~01
constrain the ML architectures to small designs. Results may improve with
multi-layer or multi-dataset training.

\textbf{No thermocouple for PODFAM.} The PODFAM dataset contains no thermocouple
reference channel, so ML methods cannot be retrained or validated on real
industrial pyrometer data. Only classical methods are applicable without
additional instrumentation.

\textbf{Fixed compression parameters.} The wavelet coefficient retention fraction
(10\,\%) and PCA component count~(3) were chosen by engineering judgement rather
than systematic optimisation.

\textbf{No real-time benchmarking.} Processing time per frame has not been
measured. Confirming that each method can operate within the pyrometer data
acquisition rate is an important requirement for industrial deployment.

\section{Future Work}
\label{sec:future}

The most impactful next step would be to instrument the PODFAM system with a
thermocouple reference channel. This would enable the ML denoisers to be
retrained on real two-colour pyrometer data and evaluated against the classical
methods on a fair basis.

Training the ML denoisers on all ten NIST layers rather than Layer~01 alone is
a straightforward improvement that could be implemented immediately with the
existing codebase, increasing the training set size tenfold.

The learnable wavelet transform proposed by Michau and Fink~\cite{learnable_wavelet}
is a natural extension of the wavelet compression method. A learnable transform
adapts the basis functions to the specific signal shape during training, which
could push the compression ratio above~$5\times$ while reducing RMSE below the
current wavelet result.

Real-time processing benchmarks are needed before any method can be deployed on
an AP\&T production line. The Kalman filter is naturally suited to real-time
deployment, but confirming that the CNN and LSTM can match the pyrometer frame
rate would open the possibility of online ML denoising.

Extending the PODFAM evaluation to files from different material types and process
parameter settings would provide a more rigorous test of the generalisability of
the classical methods beyond the ten files evaluated here.

\let\cleardoublepage\clearpage
